%\title{plantila para el congreso Ciudades y Territorios Virtuales 2018}
% From http://www.uncuyo.edu.ar/congresoctv2018/
\documentclass[11pt,a4paper]{article}

\usepackage{lipsum}
\usepackage{graphicx}
\usepackage[spanish]{babel}
\selectlanguage{spanish}
\usepackage[utf8]{inputenc}
\usepackage{authblk}
\usepackage{fancyhdr}
\pagestyle{fancy}
\renewcommand{\headrulewidth}{0pt}
\renewcommand{\footrulewidth}{0pt}
\setlength\headheight{80.0pt}
\addtolength{\textheight}{-80.0pt}
%\chead{\includegraphics[width=\textwidth]{copete_ctv.jpg}}


\fancypagestyle{firststyle}
{
%\fancyh{}
\fancyhead[c]{\includegraphics[width=15cm]{copete_ctv.jpg}}
}

\fancypagestyle{allpages}
{
%\fancyh{}
\fancyhead[c]{Mendoza}
}

\begin{document}

\pagestyle{firststyle}
%\includegraphics[width=.8\textwidth]{copete_ctv.jpg}






\title{Escribe el Título Aquí} 
\date{Octubre 2018}
\author[1]{Pedro Picapiedras\thanks{Autor a contactar: ppicapiedras@internet.net}}
\author[2]{Pablo Marmol\thanks{Co-autor: pmarmol@stone-univ.edu}}

\affil[1]{Land Use Dept, University of Nowhereland}
\affil[2]{Dept of Materials, Universite du Troyes}


\maketitle


\begin{abstract}
Las siguientes líneas son un resumen de lo que le adelantas el lector respecto del trabajo que presentarás en el congreso. Debes ser breve (no más de 100 palabras). Recuerda anticipar la situación problema, la hipoétesis de trabajo, los materiales y métodos y luego con todos estos elementos describe la investigación / experimentación que has llevado a delante y permitirá comprobar o rechazar la hipótesis. Luego han un resumen y comentario de los resultados obtenidos, una interpretación de los mismos y finaliza mencionando los aportes que tu trabajo ha hallado. Finaliza con las referencias bibliográficas, fuentes de datos y links a páginas de donde te has nutrido para este trabajo.
\end{abstract} \hspace{10pt}


\hrule


\section*{Situación problema}
\pagestyle{empty}


\begin{enumerate}
\item Pienso ...
\item Luego , existo
\item Sueno ...
\item Entonces ,\textbf{ vivo}
\end{enumerate}


\subsubsection*{Revisión Bibliográfica}
\thispagestyle{fancy}

Recuerda siempre que sugerimos no sobrepasar las 10 páginas de texto.
No numeres las páginas pues esto impide que compilemos un libro con los artículos y resúmenes. \par
\lipsum[1-2]
\section*{Materiales y métodos}
\lipsum[3]
\subsubsection*{Hipótesis de trabajo}
\lipsum[4]

\begin{figure}[ht]
\includegraphics[width=\textwidth]{Diagrama2.png}
\caption{Subtituli de fig, siempre poner esto abajo}
\end{figure}

\section*{Experimentación y caso de estudio}
\lipsum[5]

\begin{table}[!h]
\caption{Subtitulo tabla va arriba siempre}
\label{T:equipos}
\begin{center}
\begin{tabular}{| c | c | c | c | c |}
\hline
\textbf{Equipo} & \multicolumn{4}{ c |}{\textbf{Orientación del Domo Geodésico}}  \\ 
\cline{2-5}
& \textbf{Norte} & \textbf{Oeste} & \textbf{Sur} & \textbf{Este} \\
\hline

Artropodos & 3 & 0 & 0 & 0 \\ \hline
Mosquitos & 5 & 9  & 14  & 20\\ \hline
Lobesia Botrana &  0 & 1 & 0 & 2\\ \hline
Moscas & 0 & 7 & 0 & 2\\ \hline
Lepidopteros & 2 & 12  & 4  & 6\\ \hline
Cynodon dactylon & 2 & 2 & 5 & 2\\ \hline
Vinchuca & 0 & 8  & 0  & 0\\ \hline
Eficiencia Térmica & 0.6 & 0.1  & 0.3  & 0.4\\ \hline
Condición de confort (\%humedad) & 0.38 & 0.19  & 0.7  & 0.33\\ \hline

\end{tabular}
\end{center}
\end{table}
\lipsum[5]

\section*{Interpretación Aporte y Discusión}
\lipsum[5-6]

\begin{figure}[ht]
\includegraphics[width=\textwidth]{CASAS.png}
\caption{Subtituli de fig, siempre poner esto abajo}
\end{figure}
\lipsum[5-6]

\section*{Referencias Bibliográficas}

@article{ahu61,
   author={Arrow, Kenneth J. and Leonid Hurwicz and Hirofumi Uzawa},
   title={Constraint qualifications in maximization problems},
   journal={Naval Research Logistics Quarterly},
   volume={8},
   year=1961,
   pages={175-191}
}

@book{ab94,
   author* = {Charalambos D. Aliprantis and Kim C. Border},
   year = {1994},
   title = {Infinite Dimensional Analysis},
   publisher = {Springer},
   address = {Berlin}
}
*A book must have either an author field or an editor field. 
@incollection{m85,
   author={Maskin, Eric S.},
   year={1985},
   title={The theory of implementation in {N}ash equilibrium: a survey},
   booktitle={Social Goals and Social Organization},
   editor={Leonid Hurwicz and David Schmeidler and Hugo Sonnenschein},
   pages={173-204},
   publisher={Cambridge University Press},
   address={Cambridge}
}

@inproceedings{ah2006,
   author={Aggarwal, Gagan and Hartline, Jason D.},
   year={2006},
   title={Knapsack auctions},
   booktitle={Proceedings of the 17th Annual ACM-SIAM Symposium on Discrete Algorithms},
   pages={1083-1092},
   publisher={Association for Computing Machinery},
   address={New York}
}

@techreport{arrow48,
  author = {Arrow, Kenneth J.},
  title = {The possibility of a universal social welfare function},
  institution = {RAND Corporation},
  year = {1948},
  number = {P-41},
  type = {Report}
}

@unpublished{FudenbergKreps1988,
  title = {A theory of learning, experimentation, and equilibrium in games},
  author = {Fudenberg, Drew and Kreps, David M.},
  year = {1988},
  note = {Unpublished paper}
}


\end{document}